\chapter{Metodología}
\label{chap:metodologia}

El desarrollo de un sistema de inteligencia artificial aplicado a un dominio clínico
requiere no solo decisiones técnicas sobre arquitecturas y algoritmos, sino también
un marco metodológico que guíe el proceso desde la identificación del problema hasta
el despliegue de una herramienta utilizable por los profesionales. Este capítulo
describe la metodología adoptada en este trabajo, organizada en dos grandes bloques:
la metodología de gestión del proyecto y el ciclo de vida del sistema de inteligencia
artificial. A continuación se concretan los recursos hardware y software necesarios
y se describe la arquitectura general del sistema propuesto.


%% ============================================================
\section{Metodología de gestión del proyecto}
\label{sec:metodologia_gestion}
%% ============================================================

La gestión de un proyecto de \ac{IA} presenta diferencias sustanciales respecto a la
de un proyecto de tecnología de la información convencional, derivadas del ciclo de
vida específico de los sistemas de \ac{IA}, de sus riesgos particulares y de las
implicaciones morales que conlleva el procesamiento de datos sensibles. Para
estructurar la gestión de este trabajo se adopta como referencia el marco propuesto por
Miller~\cite{miller2025framework}, que entrelaza el estándar de gestión de proyectos
\emph{ISO 21502:2020-12} con el ciclo de vida de un sistema de \ac{IA} y con un
conjunto de factores de éxito orientados a evitar daños, pérdidas y perjuicios. Este
marco resulta especialmente pertinente en el presente trabajo porque los datos
manejados corresponden a menores con discapacidad, lo que sitúa la dimensión ética y
de protección de datos en el centro de todas las decisiones de diseño.

La naturaleza exploratoria de este trabajo, combinada con la restricción de operar
sobre un \emph{corpus} reducido de sesiones anotadas, aconseja un enfoque de desarrollo
iterativo e incremental dentro de ese marco. Cada iteración tiene como salida un
artefacto funcional y evaluable que puede ser revisado por el tutor y contrastado con
las anotaciones expertas antes de avanzar a la siguiente fase. Este modelo evita
comprometer recursos en el diseño de módulos completos antes de validar los supuestos
sobre los que se apoyan, lo que resulta especialmente relevante cuando los datos
disponibles son escasos y el dominio (lenguaje asistido por \ac{SAAC} en español)
carece de herramientas previas con las que comparar.

El desarrollo se organiza en fases que corresponden directamente a los objetivos
específicos definidos en el capítulo anterior: la extracción de enunciados (OE1), el
procesamiento lingüístico (OE2 y OE3), el pre-rellenado del protocolo (OE4) y la
validación (OE5). Cada fase concluye con la evaluación cuantitativa del módulo
desarrollado frente a las anotaciones del experto, lo que permite decidir de forma
fundamentada si el módulo es suficiente para avanzar o requiere refinamiento adicional
antes de integrarse en el \emph{pipeline} completo. Las tareas, plazos y dependencias
entre módulos se han gestionado mediante un tablero de seguimiento con revisiones
periódicas con el tutor.


%% ============================================================
\section{Ciclo de vida del sistema de inteligencia artificial}
\label{sec:ciclo_vida_ia}
%% ============================================================

Los marcos de referencia internacionales para el desarrollo responsable de sistemas de
\ac{IA} coinciden en estructurar el ciclo de vida en fases sucesivas que van desde la
definición del problema hasta la operación continuada del sistema y sus consecuencias.
Este ciclo de vida se divide en cuatro grandes etapas~\cite{miller2025framework}:
planificación (\emph{Plan}), desarrollo (\emph{Development}), uso (\emph{Usage}) y
consecuencias (\emph{Consequences}), subrayando que las repercusiones del sistema
sobre las personas y el entorno se prolongan más allá de la finalización del proyecto
de desarrollo. El \emph{Código de buenas prácticas para la \ac{IA} de uso general} de
la Comisión Europea~\cite{eu_ai_code_practice} y el perfil de riesgo para \ac{IA}
generativa del Instituto Nacional de Estándares y Tecnología de los Estados Unidos
(\ac{NIST} AI 600-1)~\cite{nist_ai_600_1} complementan ese ciclo con directrices
concretas sobre gobernanza, transparencia y mitigación de riesgos. Aunque este trabajo
constituye un prototipo de investigación y no un sistema en producción regulado,
adoptar estas directrices desde el inicio establece las bases para una eventual
transferencia clínica responsable.

A partir de estas referencias, el ciclo de vida del sistema desarrollado en este
\ac{TFM} se concreta en cinco fases. La fase de \emph{Definición del caso de uso} se
corresponde con la etapa de planificación; las fases de \emph{Selección de modelos} y
\emph{Diseño del sistema} desarrollan la etapa de desarrollo; la fase de
\emph{Evaluación} materializa la verificación y validación; y la fase de
\emph{Despliegue y operación} cubre las etapas de uso y consecuencias del marco de
Miller. Cada una de ellas se describe a continuación.


\subsection{Definición del caso de uso}
\label{subsec:caso_uso}

El primer paso es delimitar con precisión qué problema resuelve el sistema, para quién
y en qué condiciones de uso. En este trabajo, el caso de uso es {la asistencia al
profesional de logopedia en la cumplimentación del protocolo \ac{PAALSS}} a partir de
la grabación en vídeo de una sesión de comunicación asistida. El sistema actúa
exclusivamente como herramienta de pre-anotación: propone valores para las métricas del
protocolo, pero la decisión final sobre cada campo recae siempre en el especialista
que valida y corrige el resultado. Esta definición establece también el límite de la
responsabilidad del sistema y excluye explícitamente cualquier función de diagnóstico.

El perfil del usuario final son los profesionales de \ac{SAAC} (logopedas,
investigadores) con formación clínica pero sin necesidad de conocimientos de
programación para operar la herramienta. La restricción de procesamiento local, impuesta
por el \ac{RGPD} dado que los datos corresponden a menores con discapacidad, forma parte
de la definición del caso de uso y condiciona de forma transversal toda la arquitectura
técnica posterior.


\subsection{Selección de modelos}
\label{subsec:seleccion_modelos}

La selección de los componentes de \ac{IA} del sistema responde a tres criterios
simultáneos: adecuación técnica a la tarea, compatibilidad con el procesamiento
completamente local y disponibilidad de versiones en español o con soporte multilingüe
de calidad suficiente.

Para la transcripción automática del habla se selecciona \emph{Whisper} en su variante
\texttt{large-v3}~\cite{whisper_radford2023}, ejecutado localmente mediante la
implementación \texttt{faster-whisper}, que reduce el consumo de memoria a través de
cuantización en punto flotante de 16 bits sin pérdida apreciable de exactitud. Para la
diarización de hablantes se emplea \emph{pyannote.audio} en su configuración
\emph{pipeline} preentrenada~\cite{pyannote_plaquet2023}, complementada con el enfoque
\textbf{DiCoW}~\cite{dicow_diarization} para la transcripción condicionada por hablante
en escenarios de dos participantes.

Para el procesamiento morfosintáctico del español se emplea \emph{spaCy} con el modelo
\texttt{es\_core\_news\_lg}~\cite{spoken_spanish_pos} y \emph{Stanza}, combinados con
el modelo \emph{BETO}~\cite{beto_morphological} para las categorías morfológicas
ambiguas mediante aprendizaje con pocos ejemplos (\ac{PET})~\cite{fewshot_clinical_nlp}.
Para el módulo de pre-rellenado del protocolo se adopta una arquitectura de generación
aumentada por recuperación (\ac{RAG}) que combina una base de datos vectorial
\emph{FAISS} con un modelo de lenguaje local ejecutado mediante
\emph{Ollama}~\cite{ollama}.


\subsection{Diseño del sistema}
\label{subsec:design_sistema}

El diseño del sistema integra los modelos seleccionados en un \emph{pipeline} modular
y secuencial. La modularidad es un requisito de diseño explícito: permite sustituir
o mejorar individualmente cada componente sin rediseñar el sistema completo, y facilita
la evaluación aislada de cada módulo frente al criterio experto antes de la integración.

La arquitectura general se organiza en cuatro capas funcionales. La primera capa
realiza la extracción de enunciados a partir del vídeo de la sesión: extrae el audio,
lo transcribe con \emph{Whisper} y combina la transcripción con la diarización de
hablantes para producir una lista de Turnos Comunicativos atribuidos. La segunda capa
aplica el procesamiento lingüístico sobre dichos Turnos: etiquetado \ac{POS}, análisis
morfológico y detección de patrones sintácticos \ac{SVO}. La tercera capa calcula las
métricas del protocolo \ac{PAALSS} a partir de los enunciados procesados: diversidad
léxica, perfil morfológico, \ac{LME} y complejidad sintáctica. La cuarta capa aplica
el módulo \ac{RAG} para traducir estas métricas en el llenado estructurado de los
campos del formulario, fundamentando cada valor propuesto en los enunciados que lo
originan.


\subsection{Evaluación del sistema}
\label{subsec:evaluacion}

La evaluación sigue el protocolo de validación definido en el objetivo específico OE5.
Cada módulo se evalúa por separado antes de su integración: la transcripción mediante
la tasa de error de palabra (\ac{WER}) respecto a la transcripción manual de referencia;
la diarización mediante \emph{Diarization Error Rate} (\ac{DER}); la clasificación
morfológica y sintáctica mediante precisión, exhaustividad y F1 calculadas por
categoría; y el pre-rellenado mediante el índice \emph{kappa de Cohen} calculado campo
a campo respecto a las anotaciones del experto. Esta estrategia de evaluación modular
permite identificar los cuellos de botella del sistema y orientar los esfuerzos de
mejora de forma fundamentada.

El \emph{corpus} de evaluación son las dos sesiones anotadas manualmente por la especialista
de referencia (sesión 1: «¿Dónde están las hadas y la luna?»; sesión 2: «El ladrón de
hojas»). El reducido tamaño del \emph{corpus}, consecuencia directa de la escasez de
anotaciones expertas en este dominio, limita la generalización de los resultados y se
discute explícitamente en el capítulo de conclusiones.


\subsection{Despliegue y operación}
\label{subsec:despliegue}

El sistema se despliega como una aplicación de escritorio local que no requiere
conexión a Internet. El profesional proporciona el vídeo de la sesión a través de una 
interfaz gráfica sencilla y el sistema devuelve el formulario \ac{PAALSS} parcialmente 
cumplimentado en un formato editable. Este producto final es
un requisito del trabajo, ya que debe ofrecer una experiencia de usuario
suficiente para que un clínico sin formación en programación pueda utilizarlo de forma
autónoma.

Para validar la utilidad del prototipo más allá de las métricas cuantitativas, se
plantea la realización de una demostración en dos fases. En una primera fase interna,
se presenta el prototipo al tutor y a colaboradores del proyecto para identificar
errores de funcionamiento y aspectos de la interfaz que dificulten el uso. En una
segunda fase, se presenta la demostración a profesionales de la \emph{Fábrica de
Palabras} (asociación colaboradora que proporcionó los vídeos de la sesión), quienes
valorarán la utilidad clínica percibida del pre-rellenado, la claridad de las
propuestas del sistema y la facilidad de corrección de los campos incorrectos.


%% ============================================================
\section{Recursos hardware y software}
\label{sec:recursos}
%% ============================================================

El procesamiento completamente local impone restricciones sobre los recursos hardware
necesarios. Esta sección describe el entorno en el que se ha desarrollado y validado
el sistema.


\subsection{Recursos hardware}
\label{subsec:hardware}

El equipo de desarrollo es un ordenador personal con procesador \emph{Intel Core i7-11800H}
(8 núcleos físicos, 2,3~GHz base), 32~GB de memoria RAM y tarjeta gráfica \textbf{NVIDIA
GeForce RTX 3060} con 6~GB de VRAM. Esta configuración permite ejecutar el modelo
\emph{Whisper} \texttt{large-v3} en \ac{GPU} con cuantización \texttt{float16},
reduciendo el tiempo de transcripción de una sesión de 30 minutos a aproximadamente
4--6 minutos. El modelo de lenguaje local empleado para el pre-rellenado \ac{RAG}
(familia Llama 3 o equivalente a través de \emph{Ollama}) requiere como mínimo 8~GB de
VRAM para ejecutarse en \ac{GPU}; con la configuración disponible se pueden utilizar
modelos de 7~B de parámetros en cuantización de 4 bits. Los modelos de \ac{PLN}
(\emph{spaCy}, \emph{Stanza}, \emph{BETO}) se ejecutan en \ac{CPU} y no imponen
requisitos especiales de memoria gráfica.


\subsection{Entorno software}
\label{subsec:software}

El entorno de desarrollo se basa en Python 3.11 con gestión de dependencias mediante
\texttt{pip} y entornos virtuales. La Tabla~\ref{tab:dependencias} resume las
principales bibliotecas empleadas en cada módulo del sistema.

\begin{table}[h]
  \footnotesize
  \renewcommand{\arraystretch}{1.3}
  \setlength{\tabcolsep}{4pt}
  \begin{center}
  \begin{tabular}{|L{3.5cm}|L{4.0cm}|L{6.5cm}|}
  \hline
  \rowcolor{tema!10}
  \bft{Módulo} & \bft{Biblioteca / herramienta} & \bft{Función principal} \\ \hline
  Extracción de audio    & \texttt{ffmpeg}                  & Conversión de vídeo a WAV 16~kHz mono \\ \hline
  Transcripción \ac{ASR} & \texttt{faster-whisper}          & Transcripción local con Whisper large-v3 \\ \hline
  Diarización            & \texttt{pyannote.audio}          & Separación de hablantes \\ \hline
  \ac{PLN} español       & \texttt{spaCy}, \texttt{Stanza}  & Etiquetado \ac{POS} y morfología \\ \hline
  \ac{PLN} few-shot      & \texttt{transformers} + BETO     & Clasificación morfológica con \ac{PET} \\ \hline
  Vocabulario \ac{SAAC}  & \texttt{rapidfuzz}               & Corrección fonética y búsqueda difusa \\ \hline
  Base vectorial         & \texttt{faiss-cpu}               & Indexación de \emph{embeddings} para \ac{RAG} \\ \hline
  \ac{LLM} local         & \texttt{ollama} (SDK Python)     & Generación de texto para pre-rellenado \\ \hline
  Análisis de datos      & \texttt{pandas}, \texttt{numpy}  & Cálculo de métricas \\ \hline
  Interfaz gráfica       & \texttt{streamlit}               & \emph{Demo} de escritorio local \\ \hline
  \end{tabular}
  \end{center}
  \caption[Dependencias software]{Principales dependencias software del sistema.}
  \label{tab:dependencias}
  \renewcommand{\arraystretch}{1}
\end{table}

El sistema operativo del entorno de desarrollo es Windows 11 con acceso a \ac{GPU}
mediante CUDA 12.x.


%% ============================================================
\section{Descripción del \emph{pipeline} de procesamiento}
\label{sec:pipeline}
%% ============================================================

Esta sección detalla el diseño técnico de cada fase del \emph{pipeline}, en
correspondencia con los objetivos específicos definidos en el capítulo anterior.


\subsection{OE1: Extracción automática de enunciados}
\label{subsec:oe1_pipeline}

La extracción de Turnos Comunicativos a partir del vídeo de la sesión se lleva a cabo
en tres etapas secuenciales. En primer lugar, \texttt{ffmpeg} extrae el canal de audio
del vídeo y lo convierte a WAV mono de 16~kHz, que es el formato esperado por
\emph{Whisper}~\cite{whisper_radford2023}. En segundo lugar, \emph{Whisper} transcribe
el audio completo de la sesión y produce una lista de segmentos con su transcripción y
sus marcas temporales de inicio y fin. En esta fase se emplea un \emph{prompt} inicial
construido dinámicamente a partir del vocabulario del tablero de comunicación de la
usuaria, con el objetivo de reducir los errores fonéticos de \emph{Whisper} sobre
palabras propias del lenguaje asistido. En tercer lugar,
\emph{pyannote.audio}~\cite{pyannote_plaquet2023} diariza el audio y asigna a cada
segmento de tiempo una etiqueta de hablante (usuaria del \ac{SAAC} o interlocutor). La
fusión de la transcripción de \emph{Whisper} con la diarización de \emph{pyannote}
produce la lista de Turnos Comunicativos atribuidos.

Un reto específico de este dominio es la recuperación de los enunciados del comunicador
electrónico: la voz sintetizada que emite el dispositivo es asignada erróneamente por
el diarizador al clúster del interlocutor adulto, dado que ambas voces son fluidas. 
Para mitigar este efecto se ha diseñado un módulo de recuperación por eco
que detecta segmentos candidatos a síntesis de voz del comunicador (mediante una
puntuación ponderada de duración, número de palabras y coincidencia con el vocabulario
\ac{ARASAAC}) y los reasigna a la usuaria cuando aparecen seguidos de un eco del
interlocutor en la ventana de 8~segundos posterior.


\subsection{OE2 y OE3: Análisis lingüístico}
\label{subsec:oe2_oe3_pipeline}

Una vez obtenidos los Turnos Comunicativos de la usuaria, el módulo de procesamiento
lingüístico aplica sobre cada enunciado las siguientes operaciones. En primer lugar,
\emph{spaCy} con el modelo \texttt{es\_core\_news\_lg} etiqueta cada \emph{token} con
su categoría gramatical (\ac{POS}) y su lema, y realiza el análisis de dependencias
sintácticas. En segundo lugar, para las categorías morfológicas que presentan alta
ambigüedad sobre lenguaje asistido, como la distinción entre presente y pasado en
verbos en forma de cita, se aplica \ac{PET} sobre \emph{BETO}: el modelo recibe una
plantilla de texto que contextualiza la categoría a predecir y devuelve la etiqueta
morfológica más probable a partir de entre 20 y 50 ejemplos anotados por la experta.
En tercer lugar, el módulo de sintaxis detecta la presencia del patrón \ac{SVO} en
cada enunciado mediante heurísticos sobre el árbol de dependencias.


\subsection{OE4: Pre-rellenado del protocolo}
\label{subsec:oe4_pipeline}

El módulo de pre-rellenado traduce los enunciados procesados y las métricas calculadas
en propuestas de valores para los campos del formulario \ac{PAALSS}. La \emph{arquitectura
\ac{RAG}} recupera, para cada campo del formulario, los enunciados más relevantes de la
sesión y los suministra como contexto a un modelo de lenguaje local ejecutado mediante
\emph{Ollama}~\cite{ollama}. El modelo genera la propuesta de valor para el campo con
la instrucción explícita de basar la respuesta exclusivamente en los enunciados
recuperados y de indicar la ausencia de información cuando el contexto no sea
suficiente. La base de datos vectorial \emph{FAISS} almacena los \emph{embeddings} de
los enunciados y permite la recuperación eficiente por similitud semántica.

El resultado del módulo es un documento estructurado con las propuestas de valores para
cada campo del formulario \ac{PAALSS}, las métricas calculadas y los fragmentos de
evidencia que sustentan cada propuesta. Esto ofrece trazabilidad completa al profesional
que debe revisar y validar el pre-rellenado propuesto.


% Local Variables:
%  coding: utf-8
%  mode: latex
%  mode: flyspell
%  ispell-local-dictionary: "castellano8"
% End:
